% https://cdn3.olm.vn/upload/thuvienso/4491668113-page-36-1773022932578.png
% https://cdn3.olm.vn/upload/thuvienso/4491668113-page-37-1773022932881.png

\section*{Luyện tập chung}

\begin{lesson}
    \textbf{Ví dụ 1} Giải phương trình
    \[
        \dfrac{1}{x-1}+\dfrac{2}{x^2+x+1}=\dfrac{x^2+x}{x^3-1}. \tag{1}
    \]

    \textbf{Giải}

    ĐKXĐ: $x\ne 1$. Quy đồng mẫu hai vế của phương trình:
    \[
        \dfrac{(x^2+x+1)+2(x-1)}{(x-1)(x^2+x+1)}
        =\dfrac{x^2+x}{x^3-1}
    \]
    \[
        \dfrac{x^2+3x-1}{x^3-1}=\dfrac{x^2+x}{x^3-1}.
    \]
    Suy ra $x^2+3x-1=x^2+x$ hay $2x-1=0$.

    Giải phương trình:
    \[
        2x-1=0
    \]
    \[
        2x=1
    \]
    \[
        x=\dfrac{1}{2}\quad\text{(thoả mãn ĐKXĐ)}.
    \]
    Vậy phương trình (1) có nghiệm là $x=\dfrac{1}{2}$.

    \textbf{Ví dụ 2} Giải phương trình:
    \[
        \dfrac{x}{x+3}-\dfrac{2}{x-3}=\dfrac{-2x-6}{x^2-9}. \tag{2}
    \]

    \textbf{Giải}

    ĐKXĐ: $x\ne 3$ và $x\ne -3$.

    Quy đồng mẫu hai vế của phương trình:
    \[
        \dfrac{x(x-3)-2(x+3)}{(x+3)(x-3)}
        =\dfrac{-2x-6}{x^2-9}
    \]
    \[
        \dfrac{x^2-5x-6}{x^2-9}=\dfrac{-2x-6}{x^2-9}.
    \]
    Suy ra $x^2-5x-6=-2x-6$ hay $x^2-3x=0$.

    Giải phương trình $x^2-3x=0$:
    \[
        x(x-3)=0
    \]
    \[
        x=0\text{ hoặc }x-3=0
    \]
    \[
        x=0\text{ (thoả mãn ĐKXĐ) hoặc }x=3\text{ (không thoả mãn ĐKXĐ)}.
    \]
    Vậy phương trình (2) có nghiệm là $x=0$.

    \textbf{Ví dụ 3} Cho $a<b$. Chứng minh rằng:
    \begin{enumerate}[label=\alph*)]
        \item $2a+1<2b+2$;
        \item $-2a-5>-2b-7$.
    \end{enumerate}

    \textbf{Giải}

    a) Từ $a<b$, ta có $2a<2b$. Suy ra $2a+1<2b+1$. \hfill (1)

    Vì $1<2$ nên $2b+1<2b+2$. \hfill (2)

    Theo tính chất bắc cầu, từ (1) và (2) suy ra $2a+1<2b+2$.

    b) Từ $a<b$, ta có $-2a>-2b$. Suy ra $-2a-5>-2b-5$. \hfill (3)

    Vì $-5>-7$ nên $-2b-5>-2b-7$. \hfill (4)

    Theo tính chất bắc cầu, từ (3) và (4) suy ra $-2a-5>-2b-7$.
\end{lesson}

\begin{exercises}
    \subsection{Bài tập}

    \begin{questions}[resume=SGK]
        \item Giải các phương trình sau:
        \begin{questions}
            \item $2(x+1)=(5x-1)(x+1)$;

            \answerline
            \answerline
            \item $(-4x+3)x=(2x+5)x$.

            \answerline
            \answerline
        \end{questions}

        \item Để loại bỏ $x\%$ một loại tảo độc khỏi một hồ nước, người ta ước tính chi phí cần bỏ ra là
        \[
            C(x)=\dfrac{50x}{100-x}\quad\text{(triệu đồng)},\quad 0\le x<100.
        \]
        Nếu bỏ ra 450 triệu đồng, người ta có thể loại bỏ được bao nhiêu phần trăm loại tảo độc đó?

        \answerline
        \answerline
        \answerline

        \item Giải các phương trình sau:
        \begin{questions}
            \item $\dfrac{1}{x+2}-\dfrac{2}{x^2-2x+4}=\dfrac{x-4}{x^3+8}$;

            \answerline
            \answerline
            \answerline
            \item $\dfrac{2x}{x-4}+\dfrac{3}{x+4}=\dfrac{x-12}{x^2-16}$.

            \answerline
            \answerline
            \answerline
        \end{questions}

        \item Cho $a>b$, chứng minh rằng:
        \begin{questions}
            \item $4a+4>4b+3$;

            \answerline
            \answerline
            \item $1-3a<3-3b$.

            \answerline
            \answerline
        \end{questions}
    \end{questions}
\end{exercises}
